\section*{Preface}
\addcontentsline{toc}{section}{Preface}
\vspace{0.5cm}

\hspace{2em}These notes present a systematic summary of the material developed in Physics 131, 
Electromagnetic Theory I, with David J. Griffiths' \textit{Introduction to Electrodynamics}, 
4th Edition, serving as the principal reference. They are not intended as a verbatim transcription
of lectures--though the author does indeed copy the text verbatim in some sections--or as a replacement
for the primary text, but rather as a consolidated reference in which the principal definitions, mathematical 
methods, physical arguments, and results have been selected and reorganized for review. The presentation
also incorporates relevant material from other parts of the undergraduate physics and mathematics curriculum,
particularly vector calculus and mathematical methods encountered in Physics 116, where such material is useful
for establishing the mathematical framework required for electromagnetism. Relevant graphs were made using Python and \texttt{matplotlib},
with additional material such as relevant problem sets added.

The organization follows the development of the subject from the mathematical foundations of vector calculus
through electrostatics and magnetostatics, with emphasis on the connection between the mathematical description 
of fields and their physical interpretation. The material begins with vectors, differential and integral calculus,
curvilinear coordinates, and the fundamental theorems of vector calculus before developing Coulomb's law, electric
fields, magnetic fields, and other associated topics. These notes were prepared 
primarily as a personal reference for the author's own study and review, although they have been 
organized sufficiently to serve as a supplementary reference for classmates, peers, and other 
students who may find them useful. 

These lecture notes was prepared by the author using \LaTeXe.


\begin{flushright}
\textbf{Vercil Juan}\\
\textrm{V - BS Physics}\\
\textsc{University of the Philippines Diliman}
\end{flushright}

\vfill
\begin{center}
{\large\textsc{Physica Vincit Omnia}}
\end{center}
